% Hafta 4 — SEI CERT kuralının anatomisi
% Derleme: py -3.12 tools/latex/derle.py docs/week-4/assets/h04-15-cert-kural-anatomisi.tex
\documentclass[tikz,border=8pt]{standalone}
\usepackage{cen429-cizim}
\begin{document}
\begin{tikzpicture}[node distance=10mm and 10mm]
  \node[baslik] (b) at (0,0) {SEI CERT kuralının anatomisi};
  \node[altbaslik] at (b.south west) {her kural aynı dört parçayla gelir};

  \node[morkutu, below=16mm of b.south west, anchor=north west] (kural) {%
    \kb{Kural kimliği}\\
    INT30-C\\
    alan + sıra};
  \node[kotukutu, right=of kural] (hatali) {%
    \kb{Hatalı kod}\\
    standardın verdiği\\
    gerçek hata örneği};
  \node[iyikutu, right=of hatali] (uyumlu) {%
    \kb{Uyumlu kod}\\
    aynı işi güvenle\\
    yapan sürüm};
  \node[uyarikutu, right=of uyumlu] (risk) {%
    \kb{Risk değerlendirmesi}\\
    ciddiyet · olasılık\\
    düzeltme maliyeti};

  \draw[ok] (kural) -- (hatali);
  \draw[ok] (hatali) -- (uyumlu);
  \draw[ok] (uyumlu) -- (risk);

  \node[fit=(kural)(risk), inner sep=0pt] (grup) {};
  \node[dip, text width=155mm, below=10mm of grup.south, anchor=north] {Kuralı ezberlemeyin: hatalı ile uyumlu kodun FARKINI okuyun; öğretici olan orasıdır.};
\end{tikzpicture}
\end{document}
